% !TEX root = ../dissertationOG.tex

\clearpage % clear the prior chapter's page
\chapter{Render-Guided Volumetric Morphing with Differentiable Physics}\label{CH4_PhysMorphGS}

\noindent\textit{This chapter extends preliminary work that was presented as a poster at the ACM SIGGRAPH/Eurographics
Symposium on Computer Animation (SCA 2026); a preprint is available as \citet{song2025physmorphgs}.}
% Packages used by this chapter: amsmath, graphicx, booktabs, array, tikz, algorithm + algpseudocode, natbib (\citep).
% The details of this chapter (constants, per-step updates, term formulas, selection rule, multi-body contact,
% implementation, measures, long tables) are in the appendix "Details of Chapter 4": docs/Chapter4_appendix.tex here,
% Sections/Appendix.tex in the dissertation build. Its labels (app:ch4*, ch4:sec_below_grid, ch4:sec_cleanup,
% ch4:sec_scene, ch4:tab_b1, ch4:tab_ablation, ch4:tab_methods, ...) are referenced from this chapter.
% Dissertation build: \usepackage{tikz} and \usetikzlibrary{positioning,arrows.meta,fit,calc} are loaded in the
% preamble of VanderbiltDissertationTemplate/dissertationOG.tex.
% Dissertation build: extra stretch for this chapter only (reset at its end), so that no line runs into the margin.
\setlength{\emergencystretch}{1.5em}

% ============================================================================
\section{Overview}
\label{ch4:sec_overview}
% ============================================================================

Chapter~\ref{CH3_SHAPEMORPHING} morphed one shape into another by optimizing per-particle control deformation
gradients in a differentiable Moving Least Squares Material Point Method (MLS-MPM) simulation. The control enters the material only as an internal stress, so the
morph conserves mass and momentum and can change topology. Its objective, however, was a mass loss on the background
grid, and its appearance was produced afterwards, when the particles were meshed and smoothed for rendering. Nothing a
viewer sees, neither the silhouette of the body nor the shading of its surface, took part in the control.

In this chapter we make the picture part of the control. We optimize the same kind of physical body so that it also
\emph{looks} like the target: differentiable rendering terms, evaluated on a surface built from the body's own particles,
steer the simulation through its controls, while the controls still enter the material only as a stress and as a
rigid-free offset of its surface. Three observations shape the method. A control that is scored while it is still pushing
can win a state that springs back once it lets go. A picture asks for changes at the scale of a pixel, below the MPM
cell, which a control acting only through the grid averages away over the many particles of a cell. And the pixels of a
rendered particle cloud are not those of the surface a viewer sees. Our pipeline answers each of them:
\begin{itemize}
  \item \textbf{Settled windows.} The body is optimized window by window, and each window is scored only after the
        controls have been released and the body has settled (Section~\ref{ch4:sec_method}).
  \item \textbf{Two controls on two scales.} A capped stress increment $\delta\mathbf{F}_c$ moves the bulk through the
        grid, as in Chapter~\ref{CH3_SHAPEMORPHING}; a rigid-free normal offset $u$ of the outer particle layer shapes the
        surface below the grid (Section~\ref{ch4:sec_controls}).
  \item \textbf{A volume-conserving body} whose particles carry the volume their motion makes
        (Section~\ref{ch4:sec_body}).
  \item \textbf{Render terms on the body's exterior}, the surface of the particles' own field, weighted against a
        transport-based physics objective (Section~\ref{ch4:sec_objective}).
\end{itemize}

The chapter answers five questions in turn. \emph{What does $\delta\mathbf{F}_c$ do?} It is the bulk control: a capped
stress increment that acts through the grid, cannot change the momentum, and moves material at the scale of a cell, so
it carries the transport (Section~\ref{ch4:sec_dfc}). \emph{What does $u$ do?} It is the surface control: a rigid-free
offset of the outer particle layer along its normals, below the grid, that shapes relief finer than a cell
(Section~\ref{ch4:sec_u}). \emph{What does the render loss ask for?} What a viewer sees: the body's own exterior, drawn
from 18 views and compared with the target pictures in silhouette and shading (Section~\ref{ch4:sec_exterior}).
\emph{How does it reach the controls?} Through the surface discs to the outer particles, then directly to $u$ below the
grid, and to $\delta\mathbf{F}_c$ through the adjoint of the driven steps (Section~\ref{ch4:sec_gradpath}). \emph{And what
does it change?} Under the default rule of its weight, it redirects material while the body arrives and then polishes the
rims of thin parts, so the end is closer to the exact target by fractions of a particle spacing. Given its authority after
arrival, it writes what only the pictures carry. It carves lettering that exists only in the target pictures into a plaque
whose 3D target is smooth, and the lettering is in the material: the plaque prints it into a soft slab. It also moves
coloured material to where the pictures paint it, a layout the physics objective cannot see (Section~\ref{ch4:sec_render}).
That authority costs the body's rest, which Section~\ref{ch4:sec_costs} takes up. We measure each part against a
$\lambda=0$ twin, the same run with the render weight set to zero, and then show the method on a wider set of targets and
against other methods (Section~\ref{ch4:sec_examples}). The constants, the per-step updates and the longer tables are
collected in Appendix~\ref{app:ch4}.

\begin{figure*}[t]
    \centering
    \resizebox{\textwidth}{!}{\input{figs/ch4/pipeline_overview.tikz}}
    \caption[{Pipeline overview}]{\textbf{Pipeline overview.} Each window simulates driven and released steps, scores the released end by the
    physics terms and by the silhouette and shading of the body's exterior from 18 views, and updates both controls
    through one adjoint (the two cleanup terms of Appendix~\ref{ch4:sec_cleanup} are omitted).}
    \label{fig:overview}
\end{figure*}

% ============================================================================
\section{The Simulated Body}
\label{ch4:sec_body}
% ============================================================================

We build on the MLS-MPM simulator of Chapter~\ref{CH3_SHAPEMORPHING}, with cubic B-spline weights, Affine Particle-in-Cell (APIC)
transfers~\citep{2015-APIC,jiang2016material} and the fixed corotated elastic response~\citep{stomakhin2012energetically}. The source (a sphere) and the target mesh
are each sampled into $N$ particles of equal mass, with the source rescaled to the target's volume, and we call
$a=(V/N)^{1/3}$, the edge of the volume per particle, the \emph{sampling pitch}. The MPM cell is set by the shape, not by
$N$, so at 300k particles a cell holds about 183 of them. This is why so much of what a picture asks for lies below the
grid. In a driven step the control enters the elastic deformation,
\begin{equation}
\label{eq:control_F}
\mathbf{F}^{\mathrm{E}}_p=\big(\mathbf{F}_p+\delta\mathbf{F}_{c,p}\big)\big(\mathbf{F}^{\mathrm{P}}_p\big)^{-1},
\qquad
\boldsymbol{\tau}_p=\mathbf{P}\big(\mathbf{F}^{\mathrm{E}}_p\big)\,\mathbf{F}^{\mathrm{E}\top}_p ,
\end{equation}
and the Kirchhoff stress $\boldsymbol{\tau}_p$ enters the particle-to-grid transfer as in Chapter~\ref{CH3_SHAPEMORPHING}.

Three additions keep the body's volume and its surface in order (Appendix~\ref{app:ch4_body}). Each particle carries the
volume its motion makes, so that the smoothing of the deformation gradient used in Chapter~\ref{CH3_SHAPEMORPHING} no
longer lets thin parts in transit dilate. At every commit, a plastic assimilation turns half of the elastic stretch,
volume included, into rest shape, so the end of a morph is the body's own rest. And three position updates act below the
grid: a minimum spacing that pushes too-close particles apart with equal and opposite pushes, a rigid-free relaxation of
the outer layer toward the target's relief, and bonds that keep detached fragments with their neighbours.

% ============================================================================
\section{Two Controls: Bulk and Surface}
\label{ch4:sec_controls}
% ============================================================================

As in Chapter~\ref{CH3_SHAPEMORPHING}, we drive the morph through the material's internal deformation rather than
through prescribed trajectories or external forces. A rendering term, however, asks for changes at two scales: the bulk
of the body has to travel to the target, which the grid can carry, and its surface has to take the target's outline and
relief, much of which lies below one MPM cell. We therefore give the optimizer two controls
(Figure~\ref{ch4:fig_controls}).

\begin{figure*}[t]
    \centering
    \includegraphics[width=\textwidth]{figs/ch4/fig_controls_ppt.png}
    \caption[{Surface proximity and the outer layer}]{\textbf{Surface proximity and the outer layer.} Schematic
    cross-sections. Left: the surface proximity term pulls each surface particle toward the target surface (dashed),
    only as finely as the particle sample resolves it. Right: the layer offset moves each outer-layer particle along its
    normal by $u_p$, below the grid and free of rigid modes; the render terms read the surface discs over this layer
    from 18 views.}
    \label{ch4:fig_controls}
\end{figure*}

\subsection{What $\delta\mathbf{F}_c$ Does: The Bulk}
\label{ch4:sec_dfc}
The stress increment is the bulk control: it moves material at the scale of an MPM cell and carries the transport.
In each of the $T$ driven steps every particle has a $3\times3$ increment $\delta\mathbf{F}^{n}_{c,p}$, the control
deformation gradient of Chapter~\ref{CH3_SHAPEMORPHING}, capped per particle and step. It enters only through
Eq.~\eqref{eq:control_F}. Because the stress enters the grid as $\boldsymbol{\tau}_p(\mathbf{x}_i-\mathbf{x}_p)$ and the
B-spline weights satisfy $\sum_i w_{ip}(\mathbf{x}_i-\mathbf{x}_p)=\mathbf{0}$, the control cannot change the total linear
momentum, and since $\boldsymbol{\tau}_p$ is symmetric, APIC keeps the angular momentum as well~\citep{jiang2017angular}.
One might instead drive the body with optimized external forces; Section~\ref{ch4:sec_force} makes that replacement and
measures what changes.

\subsection{What $u$ Does: The Surface}
\label{ch4:sec_u}
The layer offset is the surface control: it shapes relief finer than a cell, which a control acting through the grid would
average away.
Every outer-layer particle has one scalar $u_p$ per window, a displacement along its normal applied in $T$ equal parts
over the driven steps:
\begin{equation}
\label{ch4:eq_u}
\mathbf{x}_p\leftarrow\mathbf{x}_p+\Big(\tfrac{1}{T}\,g_p\,u_p+s'_p\Big)\mathbf{n}_p ,
\end{equation}
where $s'_p$ is the relaxation step of Appendix~\ref{ch4:sec_below_grid} and $g_p\in\{0,1\}$ a gate. Its adjoint is the
identity times the physics' response, so the render gradient reaches it without the grid's low-pass. Three rules keep it a
surface control rather than a hand edit. It acts only on particles whose remaining transport is within one MPM cell, so it
shapes an arriving surface rather than carrying material. It is projected free of the six rigid modes of the layer after
every step of the optimizer, so it moves the body neither along nor about any axis. And it is bounded by one spacing per
window. Whether the physics holds what $u$ moves is decided by the released half of the window, in which $u$ does not act.

% ============================================================================
\section{Objective}
\label{ch4:sec_objective}
% ============================================================================

The objective has eight terms in three groups,
\begin{equation}
\label{ch4:eq_objective}
\mathcal{L}=\underbrace{\sigma\big(\mathcal{L}_{\mathrm{OT}}+\mathcal{L}_{\mathrm{prox}}+\mathcal{L}_{\mathrm{drift}}\big)
 +\mathcal{L}_{\mathrm{rel}}}_{\mathcal{L}_{\mathrm{phys}}}
 \;+\;w_c\underbrace{\big(\mathcal{L}_{\mathrm{spray}}+\mathcal{L}_{\mathrm{near}}\big)}_{\mathcal{L}_{\mathrm{clean}}}
 \;+\;\lambda\underbrace{\big(\mathcal{L}_{\mathrm{sil}}+\mathcal{L}_{\mathrm{shade}}\big)}_{\mathcal{L}_{\mathrm{render}}} ,
\end{equation}
whose formulas are given in Appendix~\ref{app:ch4_terms}. We kept a term only if switching it off changed a measured
quantity on a 19-mesh gallery and at 300k particles. The geometry and settling terms carry no weight of their own: $\sigma$
is one constant, measured once at the source, that gives the geometry energy the gradient norm of a log-mass grid loss of
the kind used in Chapter~\ref{CH3_SHAPEMORPHING}, and so fixes the unit of the cleanup weight $w_c$. There is no
regularizer of the control or of the volume.

\subsection{What the Physics Terms Ask}
The \emph{transport} $\mathcal{L}_{\mathrm{OT}}$ is the debiased Sinkhorn divergence~\citep{feydy2019sinkhorn} between the
body's mass and the target's on a loss grid. The choice matters for the morph. The nodal mass residual of
Chapter~\ref{CH3_SHAPEMORPHING} has its gradient only where the density jumps at the surface, so it cannot see a filled
region that is merely in the wrong place; the transport's gradient is a transport field over the whole body, and it is
indifferent to topology. Because the transport blurs positions within a loss cell, the \emph{surface proximity}
$\mathcal{L}_{\mathrm{prox}}$ asks, at every outer point of the target, whether body material is there. Two settling terms
ask for an end at rest: the \emph{end drift} $\mathcal{L}_{\mathrm{drift}}$ charges the velocity left at the released end,
and the \emph{released motion} $\mathcal{L}_{\mathrm{rel}}$ charges all motion during the released half; without the
latter, runs stop while the body is still moving. The \emph{cleanup} group is optional. Its two local terms,
Eqs.~\eqref{ch4:eq_spray} and~\eqref{ch4:eq_near}, only tidy what the transport cannot resolve below a loss cell:
particles sprayed off the body, and particles in a thin band just off the target surface. We use $w_c=0.2$ throughout,
and $w_c=0$ removes the group. In the leave-one-out tests that fixed the objective, removing the spray pull left the
shape unchanged but left more isolated particles and lengthened the runs, and removing the near band left more thin
target points uncovered at 40k particles (Appendix~\ref{ch4:sec_cleanup}).

\subsection{What the Render Loss Asks For}
\label{ch4:sec_exterior}
The render loss asks for what a viewer sees, and it reads the surface of the body, not its particle cloud. The body's \emph{exterior} is the zero set of Zhu
and Bridson's field~\citep{10.1145/1073204.1073298} of the particles, covered by surface discs that follow the particles. The
\emph{silhouette} term compares the discs' coverage in 18 orthographic views with the target's, penalizing missing
coverage inside the target more than excess outside it; the \emph{shading} term compares headlight-lit images of the
discs. The target pictures are drawn by the same operator from the exterior of the target sample, averaged over eight
independent samples of the mesh so that the reference carries no single sample's noise (Appendix~\ref{app:ch4_render}).
Reading the render terms on the exterior rather than on the particle cloud is what makes them act on the surface a viewer
sees: read on the cloud, they improve the cloud's own silhouette and leave the displayed surface where the twin has it
(Section~\ref{ch4:sec_ablation}). The images in this chapter draw the same exterior with surface discs in a studio look;
every image is a committed simulated state.

% ============================================================================
\section{Morphing Method}
\label{ch4:sec_method}
% ============================================================================

Given the body, its controls and the objective, the morph is an optimization over a sequence of short simulation
networks, as in Chapter~\ref{CH3_SHAPEMORPHING}. Here each network is a settled window, and a window is committed only
after the body has been released and scored.

\subsection{Settled Windows}
A window simulates $T$ driven steps with the controls, then $T$ released steps with the controls at zero while the
physics keeps running. Every term of the objective is evaluated on the released end state and the released velocities.
Rebound, ringing and anything else that the control holds only while it pushes therefore cost the control directly. The
optimizer runs a fixed number of iterations per window; the window is then committed, and the next one starts from the
committed state (Algorithm~\ref{ch4:alg_window}).

\subsection{How the Render Reaches $\delta\mathbf{F}_c$ and $u$}
\label{ch4:sec_gradpath}
The render loss is a function of pictures, and the pictures are drawn from discs that sit on the outer particles. Its
gradient therefore lands first on the discs and, through the exterior's field, on the positions of the outer particles at
the released end (Figure~\ref{ch4:fig_gradpath}). On its way to the controls it is combined with the physics gradient: the
part that opposes the physics is removed, and the rest is scaled by the render weight $\lambda$
(Section~\ref{ch4:sec_lambda}). It then reaches the two controls by two paths. The short path goes to $u$: the offset
moves those same particles along their normals, so the gradient arrives below the grid, without the grid's low-pass, and
can carve relief finer than a cell. The long path runs backward through the released and the driven steps of the window,
through the transfers between particles and grid, to every $\delta\mathbf{F}^n_c$, and reshapes the body at the scale of a
cell. The render has no other route to the material: no external force, and no edit of positions outside the two
controls. Section~\ref{ch4:sec_where} measures how the two paths share the work.

\begin{figure*}[tbp]
    \centering
    \IfFileExists{figs/ch4/fig_gradpath.pdf}{\includegraphics[width=\textwidth]{figs/ch4/fig_gradpath.pdf}}{%
    \IfFileExists{figs/ch4/fig_gradpath.png}{\includegraphics[width=\textwidth]{figs/ch4/fig_gradpath.png}}{%
    \fbox{\parbox{0.9\textwidth}{\centering fig\_gradpath pending}}}}
    \caption[{How the render loss reaches the two controls}]{\textbf{How the render loss reaches the two controls.}
    (a) Bulk stress $\delta\mathbf{F}_c$ and normal offsets $u$ act for $T$ driven steps.
    (b) Physics continues for $T$ released steps with both controls zero; surface discs follow the released body.
    (c) Rendering compares this exterior with target images. Each loss is differentiated through the entire window,
    and its control-space gradients are combined using PCGrad and $\lambda$ before updating both controls.
    The dashed curve is the 3D placement target; image targets may supply additional content.
    Geometry and image panels are schematic; cleanup gradients are omitted.}
    \label{ch4:fig_gradpath}
\end{figure*}

\subsection{Gradients and the Render Weight}
\label{ch4:sec_lambda}
The physics, cleanup and render terms are differentiated separately through the same $2T$-step adjoint back to every
$\delta\mathbf{F}^n_{c}$ and to $u$. Let $\mathbf{g}_p$ be the gradient of the physics core and $\mathbf{g}_r$ that of the
render term. The render gradient loses its component against the physics gradient when the two
conflict, as in PCGrad~\citep{yu2020pcgrad},
\begin{equation}
\tilde{\mathbf{g}}_r=\mathbf{g}_r-\frac{\min(0,\langle\mathbf{g}_p,\mathbf{g}_r\rangle)}{\lVert\mathbf{g}_p\rVert^{2}}\,\mathbf{g}_p ,
\end{equation}
and the render weight is set at each window's first gradient by
\begin{equation}
\label{ch4:eq_lambda}
\lambda\,\lVert\tilde{\mathbf{g}}_r\rVert=\tfrac12\,\lVert\mathbf{g}_p\rVert ,
\end{equation}
so that the render term holds one third of the first gradient in every window. The step direction is
$\mathbf{g}=\mathbf{g}_p+\lambda\tilde{\mathbf{g}}_r+\mathbf{g}_{\mathrm{clean}}$.

This rule keeps the render's \emph{share} of a step constant, but not its \emph{authority}. While the body arrives, the
physics gradient falls by orders of magnitude, and $\lambda$ falls with it. The render term therefore acts strongly while
the body is on its way and only weakly once it has arrived, a consequence we measure in
Section~\ref{ch4:sec_render_effect} and take up again in Section~\ref{ch4:sec_costs}.

\subsection{A Second Mode: The Render Keeps Its Weight}
\label{ch4:sec_held}
The rule of Eq.~\eqref{ch4:eq_lambda} ties the render's weight to the physics residual. On the 300k bunny, $\lambda$ falls
from 0.32 at the first window to below 0.003 once the body has arrived, and the render's share of the outer layer's pull
settles near a fifth. The render then has no authority left exactly when there is something only it could decide: which
material ends where, and relief finer than the 3D target resolves. The physics objective cannot see either, because
every physics term reads only the set of particle positions.

In a second operating mode, the render keeps its first-window weight: $\lambda$ is calibrated once by
Eq.~\eqref{ch4:eq_lambda} and then held, or only the weight of the term that carries the picture-only content, such as a
colour term, is held while the other render terms follow the rule. The purpose of this mode is content that only the
pictures carry: relief that exists in the target pictures alone, and the placement of material by its appearance. After
arrival the fixed-weight cleanup terms are among the forces that push back against the render, so $w_c=0$ is a natural
companion of this mode, although the runs of Section~\ref{ch4:sec_write} keep it. The mode has a cost: a held weight does
not vanish once the body has arrived, so the body is still moving when the run stops, and the finest shape metrics relax
slightly. A weight that keeps the render's authority where it has something to decide and lets the body come to rest is
ongoing work (Section~\ref{ch4:sec_costs}).

\subsection{Step, Commit and Stop}
The step is Adam's~\citep{kingma2014adam} with a backtracking line search, as in
Chapter~\ref{CH3_SHAPEMORPHING}; after each trial $\delta\mathbf{F}_c$ is capped and $u$ is bounded and made rigid-free.
The moments of $u$ carry from one window to the next, so that its first step in a window does not reverse the last one.
A window's result is committed together with the plastic assimilation, rejected if it raises the objective markedly, and
the run stops when the objective no longer improves (Appendix~\ref{app:ch4_method}).

\begin{algorithm}[t]
  \caption{Settled-window morph}\label{ch4:alg_window}
  \begin{algorithmic}[1]
    \Require source and target samples, $T$ driven and $T$ released steps, $i_{\max}$ iterations per window
    \State $\mathbf{s}\gets$ the source body at rest;\ $\delta\mathbf{F}^{\mathrm{prev}}_c\gets\mathbf{0}$
    \While{the stop rule has not fired}
      \State $\delta\mathbf{F}_c\gets\tfrac12\delta\mathbf{F}^{\mathrm{prev}}_c$ if that beats zero, else $\mathbf{0}$;\ $u\gets 0$
      \For{$i=1,\dots,i_{\max}$}
        \State simulate $T$ driven steps ($\delta\mathbf{F}_c$, $u$ act) and $T$ released steps from $\mathbf{s}$
        \State evaluate $\mathcal{L}_{\mathrm{phys}}$, $\mathcal{L}_{\mathrm{clean}}$, $\mathcal{L}_{\mathrm{render}}$ on the released end
        \State $\mathbf{g}_p,\mathbf{g}_{\mathrm{clean}},\mathbf{g}_r\gets$ the adjoint through the $2T$ steps
        \State $\tilde{\mathbf{g}}_r\gets$ PCGrad$(\mathbf{g}_r,\mathbf{g}_p)$;\ at $i=1$:\ $\lambda\gets\tfrac12\lVert\mathbf{g}_p\rVert/\lVert\tilde{\mathbf{g}}_r\rVert$
        \State Adam step on $\mathbf{g}_p+\lambda\tilde{\mathbf{g}}_r+\mathbf{g}_{\mathrm{clean}}$ with a backtracking line search
        \State cap $\delta\mathbf{F}_c$;\ bound $u$ and remove its rigid modes
      \EndFor
      \State replay the accepted candidate;\ assimilate half the elastic stretch into $\mathbf{F}^{\mathrm{P}}$
      \If{the merit rose by more than $5\,\%$}
        \State reject:\ keep $\mathbf{s}$
      \Else
        \State commit:\ $\mathbf{s}\gets$ the released end state;\ $\delta\mathbf{F}^{\mathrm{prev}}_c\gets\delta\mathbf{F}_c$
      \EndIf
    \EndWhile
    \State \Return the trajectory up to the best window
  \end{algorithmic}
\end{algorithm}

\subsection{What the Body Conserves}
\label{ch4:sec_conservation}
Mass is exact by construction. The stress control cannot change the linear or angular momentum, the minimum spacing's
push is pairwise and central, and the relaxation and $u$ are free of the rigid modes of the layer. The one exception is a
weak ambient drag, which is not angular-momentum consistent, so we report rotation as a kinematic quantity. Over whole
300k morphs the centre of mass moves by less than $10^{-4}$ pitches, and the body's physical volume ends within half a
percent of the target's.

\subsection{Bodies in a Scene}
A single world-fixed objective cannot let a body fly, spin or land. For scenes with several bodies, contact and gravity,
each body deposits into its own grid field, so that bodies meet with a momentum-conserving contact rule instead of
sticking, and each body's objective is read in its own shape-matching frame, so that flight and tumbling cost nothing.
Every window is committed, and the targets move along their transport so that bodies meet half-formed. The single-body
path is unchanged by these extensions (Appendix~\ref{ch4:sec_scene}).

% ============================================================================
\section{What the Render Does}
\label{ch4:sec_render}
% ============================================================================

\subsection{Setup}
\label{ch4:sec_protocol}
Every morph of the main evaluation starts from a sphere and targets one of six meshes at $N=300{,}000$: the bunny, the
dragon, the teapot, bob, the armadillo and homer, three of them with handles. A 19-mesh gallery runs at $N=40{,}000$, and
twenty further targets and shape pairs run at 300k. We name the compared arms by what they contain: \emph{the full method}
($\delta\mathbf{F}_c$ + layer $u$ + render loss); \emph{its $\lambda=0$ twin} ($\delta\mathbf{F}_c$ + layer $u$);
\emph{$\delta\mathbf{F}_c$ only}; and \emph{the F-control of Chapter~\ref{CH3_SHAPEMORPHING}}.

No result is scored against the sample it was optimized toward. We read each end against the exact target mesh on views
the loss never sees, using standard silhouette, surface and image metrics; at the particle level against an independent
sample of the mesh, whose second sample defines the \emph{sampling floor}, a perfect placement; and through the render
terms' own operators, which locate where the render acts but, as we show below, also reward the arrangement of the outer
layer. Appendix~\ref{app:ch4_measures} defines each measure and the rule, fixed before any metric was computed, for when a
metric \emph{sees} the render.

\subsection{Where the Render Acts}
\label{ch4:sec_where}
Figure~\ref{ch4:fig_rgv} follows the render's pull on a deforming 300k bunny. It lands on the ears, the face and the feet:
early on mostly in agreement with the physics, later more often against it, where the projection removes the opposing
part. Both paths of Section~\ref{ch4:sec_gradpath} carry it. With the render weight held (Section~\ref{ch4:sec_held}) and
$u$ switched off, so that the render reaches the body only through $\delta\mathbf{F}_c$, a coin keeps about a third of the
lettering the render carves into it with $u$, while the lettering on a bunny's flank is carved as deeply as with $u$. The
surface offset carries most of the lettering on the coin, and the stress increment alone carries it on the bunny
(Appendix~\ref{app:ch4_held}).

\begin{figure*}[p]
    \centering
    \includegraphics[width=\textwidth]{figs/ch4/fig_render_gradient_early.jpg}\\[4pt]
    \includegraphics[width=\textwidth]{figs/ch4/fig_render_gradient_late.jpg}
    \caption[{The render gradient on the deforming 300k bunny, early (top) and late (bottom)}]{\textbf{The render gradient on the deforming 300k bunny, early (top) and late (bottom).} Each block shows the
    render's share of each particle's pull, its strength, its agreement with the physics, and what it changed in
    $\delta\mathbf{F}_c$.}
    \label{ch4:fig_rgv}
\end{figure*}

\subsection{Under the Default Rule: Arriving, Then Polishing}
\label{ch4:sec_render_effect}
Under the default rule the render's weight follows the physics gradient. Figure~\ref{ch4:fig_render_effect} shows the end:
against its twin, the full method ends closer to the exact target on every mesh and seed, by a small margin that sits on
the rims of thin parts. We follow the render through a run in three phases: the arrival, the time after it, and the end.

\begin{figure*}[tbp]
    \centering
    \includegraphics[width=0.9\textwidth]{figs/ch4/fig_render_effect_v2.png}
    \caption[{What the render loss changes at the end of a 300k morph}]{\textbf{What the render loss changes at the end of a
    300k morph.} (a, b) One row per mesh and seed, read against the exact target mesh on views the loss never sees: grey is
    the $\lambda=0$ twin ($\delta\mathbf{F}_c$ + layer $u$), orange the full method, and the black tick a perfect particle
    sample of the target. The outline error is lower in all 22 pairs, by a median of 16\,\%, and the surface error at half a
    pitch in all 22, by a median of 36\,\%. (c) The difference sits on the rims of thin parts: the end silhouettes of the
    bunny's ears and the dragon's tail tip (seed 97, held-out views) against the target's outline.}
    \label{ch4:fig_render_effect}
\end{figure*}

\textit{During the arrival, the render redirects material.} In the first second the full method's surface parts from its
twin's by far more than two runs of the same arm with different seeds part from each other; on the dragon the gap reaches
almost a pitch. The render carries only a small share of the surface's motion in that second, but that second holds most
of all the motion the morph makes, on the ears, the face and the feet of Figure~\ref{ch4:fig_rgv}.

\textit{After the arrival, it only polishes.} Once the body has arrived, the physics gradient has fallen by more than an
order of magnitude, and with it, by the rule of Section~\ref{ch4:sec_lambda}, the render weight and the step. The two arms'
surfaces return to a difference no larger than the seed-to-seed noise. What the render term still does is hold a sub-pitch
offset on silhouette rims and thin parts against the physics' own refit. A causal test shows this
(Figure~\ref{ch4:fig_switch}): switched off at 6\,s, the full method loses its lead within about a second; switched on at
6\,s, the twin builds it within a third of a second and ends level with the run that had the render throughout. Most of
what the render asks for lies below the MPM cell, and it is $u$, not $\delta\mathbf{F}_c$, that reaches it.

\textit{At the end, the gains are small but consistent.} Table~\ref{ch4:tab_d157} reads the end against the exact target
meshes, on views the loss never sees. At 300k every standard metric but one improves, in every mesh and in nearly every
pair of runs; the one exception is saturated. The improvement is real but small in absolute terms: the Boundary IoU~\citep{cheng2021boundaryiou} error
goes from 0.068 to 0.054, near a perfect sample's 0.052 (Figure~\ref{ch4:fig_render_effect}). The gain is the same on views far from every training view, so it
is not specific to the pictures the loss sees; Table~\ref{ch4:tab_b1} gives it per mesh. At 40k the same metrics move in
the same direction, by about half as much on the boundary metrics and less elsewhere. The particle-level measures do not
see these gains: the particle silhouette, chamfer distance, thin coverage, centre of mass and rest stay within the twin's
range (Figure~\ref{ch4:fig_render_end}), because the binary particle silhouette resolves almost two pitches per pixel while
the render changes the exterior by fractions of a pitch (Figure~\ref{ch4:fig_iou}).

\begin{table}[t]
\centering
\footnotesize
\setlength{\tabcolsep}{4pt}
\caption[{The full method against its $\lambda=0$ twin, read against the exact target mesh}]{\textbf{The full method against its $\lambda=0$ twin, read against the exact target mesh.} $r$ is the median share
of the twin's error the full method removes (for PSNR, the gain); $^\ast$ marks a metric that sees the render by the rule
fixed beforehand.}
\label{ch4:tab_d157}
\begin{tabular}{lcccc}
\toprule
 & \multicolumn{2}{c}{300k (6 meshes, 22 pairs)} & \multicolumn{2}{c}{40k (19 meshes, 57 pairs)} \\
\cmidrule(lr){2-3}\cmidrule(lr){4-5}
metric & $r$ & pairs / meshes & $r$ & pairs / meshes \\
\midrule
\multicolumn{5}{l}{\textit{silhouettes, held-out views (1024\,px)}}\\
Boundary IoU (2\,\% of the diagonal) & $+16.3\,\%^\ast$ & 22 / 6 & $+7.1\,\%^\ast$ & 52 / 17 \\
boundary F-score (DAVIS, 12\,px)      & $+0.9\,\%$       & 12 / 3 & $+12.5\,\%$      & 31 / 11 \\
contour Chamfer                       & $+20.2\,\%^\ast$ & 21 / 6 & $+8.4\,\%^\ast$  & 54 / 18 \\
contour HD95                          & $+19.6\,\%^\ast$ & 21 / 6 & $+8.1\,\%^\ast$  & 50 / 18 \\
\multicolumn{5}{l}{\textit{3D, visible exterior}}\\
F-score at 0.5 pitch                  & $+36.4\,\%^\ast$ & 22 / 6 & $+8.0\,\%^\ast$  & 46 / 16 \\
F-score at 1 pitch                    & $+14.4\,\%^\ast$ & 22 / 6 & $+3.5\,\%$       & 39 / 15 \\
normal consistency                    & $+9.2\,\%^\ast$  & 21 / 6 & $+2.8\,\%$       & 42 / 15 \\
Chamfer-L1                            & $+12.9\,\%^\ast$ & 22 / 6 & $+4.3\,\%$       & 46 / 16 \\
\multicolumn{5}{l}{\textit{images, held-out views (512\,px)}}\\
shaded PSNR                           & $+0.73$\,dB$^\ast$ & 22 / 6 & $+0.30$\,dB$^\ast$ & 53 / 18 \\
shaded SSIM                           & $+10.7\,\%^\ast$ & 22 / 6 & $+3.0\,\%$       & 52 / 18 \\
shaded LPIPS                          & $+7.0\,\%^\ast$  & 22 / 6 & $+2.5\,\%$       & 47 / 17 \\
normal-map PSNR                       & $+0.96$\,dB$^\ast$ & 22 / 6 & $+0.36$\,dB$^\ast$ & 53 / 18 \\
normal-map SSIM                       & $+9.0\,\%^\ast$  & 22 / 6 & $+2.7\,\%$       & 51 / 17 \\
normal-map LPIPS                      & $+10.0\,\%^\ast$ & 22 / 6 & $+3.8\,\%$       & 51 / 17 \\
\midrule
\multicolumn{5}{l}{\textit{for comparison}}\\
particle silhouette IoU (128\,px)     & $+2.3\,\%$       & 14 / 4 & $+4.5\,\%$       & 48 / 19 \\
mask IoU against the raw mesh         & $+5.2\,\%$       & 22 / 6 & $-0.9\,\%$       & 19 / 5 \\
\bottomrule
\end{tabular}
\end{table}

\textit{Part of the render's lead on its own term is arrangement, not shape.} At the end of the 300k runs the full
method's render term lies below that of every uniform sample of the true mesh, which no body of the exact shape, sampled
uniformly, can do. What goes below that floor is an outer layer arranged more evenly, and hence an exterior smoother, than
any random sample; on the bunny at least two fifths of the full method's lead over its twin on the render term is of this
kind, and on the dragon at least a fifth. This is why we read shape only on the exact mesh and use the render terms' own
operators, as in Figures~\ref{ch4:fig_render_end} and~\ref{ch4:fig_switch}, only to locate the render's action.

In short, under the rule, the render term reshapes the arrival and polishes the end; after arrival it decides nothing
that the physics objective leaves open, which is what holding its weight changes (Section~\ref{ch4:sec_write}). It also
has a small cost: it deforms some bodies more, and the ambient drag, which does not respect angular momentum, acts more on
a body that moves more. Appendix~\ref{app:ch4_results} gives the measurements behind this section.

\begin{figure*}[t]
    \centering
    \includegraphics[width=\textwidth]{figs/ch4/fig_render_switch.png}
    \caption[{The render term switched at 6\,s}]{\textbf{The render term switched at 6\,s.} Switched off, the full method loses its lead within about a second;
    switched on, the twin builds it within a third of a second.}
    \label{ch4:fig_switch}
\end{figure*}

\subsection{What the Render Can Write}
\label{ch4:sec_write}
Under the default rule of Section~\ref{ch4:sec_lambda} the render only polishes rims once the body has arrived, because its
weight follows the physics gradient. That is a property of the weight, not of what the render can reach. With its weight held after arrival (Section~\ref{ch4:sec_held}), the render
does what the physics objective cannot: it writes content that only the pictures carry.

\textit{Relief that exists only in the pictures.} The 3D target is a smooth plaque; the word SIGGRAPH on its front and
VANDERBILT on its back exist only in the render's target pictures (Figure~\ref{ch4:fig_letters}). The twin, which has the
layer offset $u$ but no render loss, ends smooth. With the render weight held, the full method carves both words into the
particles, to 58\,\% of the relief that a 300k sample of the lettered mesh itself carries. This also answers whether the
visible relief could come from $u$ alone: $u$ acts in both arms, and only the arm that sees the pictures writes the words.
Under the rule, the render's weight has gone by the time the plaque has formed, and the words stay at the twin's level on
a 40k plaque. Reading the physics placement terms only at the scale of the loss cell, which leaves the finer scale to the
render, lets the held render carve more of the lettering on 40k bodies, up to 0.83 of the ceiling. The cost is the one
named in Section~\ref{ch4:sec_held}: the body is still moving when the run stops, and the plaque's ends are slightly
under-filled.
%% SWAP-IN: the 300k ripples (job 55) and the 500k ripples (job 53) go here when they land (figure and one number).

\begin{figure*}[tbp]
    \centering
    \includegraphics[width=\textwidth]{figs/ch4/fig_render_letters.jpg}
    \caption[{Lettering that exists only in the target pictures}]{\textbf{Lettering that exists only in the target
    pictures.} A 300k plaque whose 3D target is smooth: with the render weight held, the full method carves SIGGRAPH
    (front) and VANDERBILT (back) into the particles, while its twin, which has the layer offset $u$, ends smooth.}
    \label{ch4:fig_letters}
\end{figure*}

\textit{Placement by appearance.} Each particle carries a colour that is fixed at the source and moves with it, and the
target pictures are painted. The physics objective is provably blind to the paint: the twin's trajectory is bit for bit
the same under any paint. In the scene of Figure~\ref{ch4:fig_colour}, the source is red on its upper half and the pictures
paint the head half of the bunny red. The twin keeps the transport's top-and-bottom split, and under the rule the full
method leaves the layout almost where the twin leaves it. With the colour term's weight held, the full method moves
material so that the colours land where the pictures put them, and recovers about 70\,\% of the layout's error. It too
ends while the colour is still moving, and its finest surface metric relaxes slightly. Appendix~\ref{app:ch4_held} gives
the measurements behind this section.
%% SWAP-IN: the colour band at 300k (job 55) and at 500k (job 54) goes here when it lands (figure and one number).

\begin{figure*}[tbp]
    \centering
    \includegraphics[width=\textwidth]{figs/ch4/fig_render_colour.jpg}
    \caption[{Placing material by its appearance}]{\textbf{Placing material by its appearance.} The physics objective
    cannot see which particle carries which colour; with the colour term's weight held, the full method moves the red
    material to the head half that the target pictures paint, while its twin keeps the transport's top-and-bottom split
    (bunny, 40k particles, front and back).}
    \label{ch4:fig_colour}
\end{figure*}

\textit{Where the render acts when it keeps its weight.} Figure~\ref{ch4:fig_held_share} repeats the share map of
Figure~\ref{ch4:fig_rgv} on the 300k bunny, under the rule and with the render weight held. Under the rule, the render's
share of the outer layer's pull settles at 0.14--0.23 after arrival; held, it is 0.78--0.90, and the render, not the
physics, then decides what the outer layer does for the rest of the run. That is how it writes the content of
Section~\ref{ch4:sec_write}, and it is also its price: the body's kinetic energy at the end is three orders of magnitude
above the rule's, and a few more thin target points are left uncovered (Appendix~\ref{app:ch4_held}).

\begin{figure}[tbp]
    \centering
    \includegraphics[width=0.62\linewidth]{figs/ch4/fig_render_held_share.jpg}
    \caption[{Where the render acts when it keeps its weight}]{\textbf{Where the render acts when it keeps its weight.}
    The render's share of each particle's pull on the 300k bunny, from grey (the physics alone) to orange (the render
    alone): under the rule it fades after arrival, held it dominates the outer layer.}
    \label{ch4:fig_held_share}
\end{figure}


\subsection{The Letters Are in the Material}
\label{ch4:sec_stamp}
A carving that lived only in the displayed surface would be a decal. To test where the lettering is, we press a finished
plaque, letters down, into a soft slab and lift it again, with the controls off (Figure~\ref{ch4:fig_stamp}). The full
method's plaque prints SIGGRAPH into the slab. The twin, given the same lettering afterwards as a decal on its displayed
surface, prints nothing: in a still picture the decal can even look crisper, but it cannot press. At 300k the print
correlates with the letters at 0.46, and at 500k it reads under raking light. This is what separates the render loss from
a physics-only morph followed by a surface fit to the pictures: the fit changes the picture, while the render loss changes
the body. Appendix~\ref{app:ch4_held} describes the press.
%% SWAP-IN: the stiffness result (job 51) goes here when it lands: whether a stiffer plaque keeps and prints the letters.

\begin{figure*}[tbp]
    \centering
    \includegraphics[width=\textwidth]{figs/ch4/fig_render_stamp.jpg}
    \caption[{The lettering is in the material}]{\textbf{The lettering is in the material.} Plaques lettered as in
    Figure~\ref{ch4:fig_letters}, pressed letters down into a soft slab and lifted. Top: at 500k under raking light, the
    full method's plaque face and the print it leaves, mirrored to read, against its twin with the lettering added
    afterwards as a decal. Bottom: the depth of the print at 300k, darker where deeper.}
    \label{ch4:fig_stamp}
\end{figure*}

% ============================================================================
\section{Examples and Comparisons}
\label{ch4:sec_examples}
% ============================================================================

\begin{figure*}[t]
    \centering
    \includegraphics[width=\textwidth]{figs/ch4/fig_morph_sequence_v2.jpg}
    \caption[{From a sphere to five targets at 300k particles}]{\textbf{From a sphere to five targets at 300k particles.} Every
    row is shown at the same times and at its end. The body rises as one piece and the features grow out of it in the first
    second or two; bob's ring opens a little later and the dragon's narrowest gaps open last. A separate piece of the display
    surface that holds fewer than ten particles is not drawn (Appendix~\ref{app:ch4_render}).}
    \label{ch4:fig_morph}
\end{figure*}

\subsection{From a Sphere to Five Targets}
Figure~\ref{ch4:fig_morph} shows five of the targets. The body rises as one piece, and the features grow out of it within
the first second or two: the bunny's ears and the armadillo's limbs are in place by 1.5\,s, bob's ring opens a little
later, and the dragon's horns and claws and the narrow gaps between its coils come last. Most of the particle movement
happens in these first seconds; what follows is a long, slow tail that fills most of the delivered frames. The ends are
close to what a perfect placement of the particles would give: the particle silhouette lies just below the sampling floor
and the chamfer distance below it, because the ends are more even than any random sample. The bodies end at rest, and no
run sheds a serious piece. Section~\ref{ch4:sec_render_effect} reads these ends against their twins.

\subsection{The 40k Gallery}
\label{ch4:sec_gallery40}
On the 19-mesh gallery at $N=40{,}000$ the render acts in the same direction as at 300k, at about half the size: the
particle silhouette is ahead of the twin's on all but one mesh by a small margin, and against the exact meshes the outline
improves on all 19. One reading goes the other way: the silhouette area against the raw mesh is slightly worse, because the
exterior leans slightly outward. The lean comes from the silhouette term's asymmetric weighting of missing against excess
coverage, which at 40k, unlike at 300k, is not balanced by an inward pull; Appendix~\ref{app:ch4_gallery} gives the
measurement and the next test.

\subsection{A Gallery of 300k Morphs}
\label{ch4:sec_showcase}
To see the method beyond the six main targets, we ran it unchanged from a sphere to the 13 other gallery meshes and to the
Happy Buddha, a target of genus six, and between six pairs of shapes; Figure~\ref{ch4:fig_showcase} shows all of them but
the three letters, A, V and C. Every run but one
reaches its target nearly as closely as the main six do, comes to rest, and sheds no serious piece. Topology changes as
the targets ask. From the teapot to bob, the teapot's handle closes within the first second and bob's ring opens later,
so the body ends with the target's single hole; from a sphere, the Happy Buddha's openings form. Where twins were run, the
render's lead has the same direction and size as on the main six.

Three examples show where the method still pays. On the letter C the run stops early while the body still flows, because
the stop rule reads the objective's progress, not rest. Between two shapes in different orientations, such as spot and the
cow, the transport carries the head over the back; yawing the source first makes the same morph direct and twice as fast.
And a less compact source such as the armadillo runs on a cell that is coarse relative to the body, and passes through
thin, perforated sheets on its way to homer.

\begin{figure*}[p]
    \centering
    \includegraphics[width=\textwidth]{figs/ch4/fig_showcase_montage_v2_interim24.jpg}
    \caption[{Twenty-four 300k morphs, each shown at its start, halfway and end}]{\textbf{Twenty-four 300k morphs, each shown
    at its start, halfway and end.} The gallery of Section~\ref{ch4:sec_showcase} without the three letters. The last panel
    uses an MPM cell set from the source's volume, an option outside the method, for comparison with the panel before it.}
    \label{ch4:fig_showcase}
\end{figure*}

\subsection{Materials and External Forces}
Every material we tried reaches its target: a soft solid arrives more slowly, a stiff one faster but less at rest, and
the render term's lead holds in each. Under gravity, the single-body window stops early, because a falling body raises the
world-fixed objective. Through the scene path, with the target standing on the floor, the bunny and the dragon reach their
targets on every seed, stand upright at the target's height and keep their horizontal momentum. They do not come to rest:
each commit's assimilation makes part of the weight's sag rest shape and the controls raise it again. The next definition
assimilates only the controls' share of each window's deformation (Appendix~\ref{app:ch4_gravity}).

\subsection{Ablations}
\label{ch4:sec_ablation}
Switching off one part at a time at 300k (Table~\ref{ch4:tab_ablation}) gives three readings. First, $u$ is the main shape
lever: without it every surface measure is several times worse and the body never comes to rest, although the early
floating particles almost vanish, because $u$ is also their source. Second, reading the render on the exterior beats reading
it on the particle cloud on every measure of the surface, while the cloud read wins on the particle measures it optimizes;
each read wins on the representation it scores, and the exterior is the one a viewer sees. Third, the volume-conserving
body, the minimum spacing and the choices made against late motion each earn their place on what they address.

\subsection{Stress Control against Force Control}
\label{ch4:sec_force}
Physics-based shape control drives matter toward a target with external
forces~\citep{treuille2003keyframe,2004-FluidControlAdjointMethod,10.1145/1015706.1015743,chen2024eigenfluids,li2025flowmapadjoint}.
To test our choice, we replaced $\delta\mathbf{F}_c$ by an optimized external force per particle and step and kept
everything else. A plain force puts momentum in at every window, which the objective then has to pull back. Projecting the
force onto zero net force and torque removes that momentum exactly, but not the shape deficit: both force arms end with
several times the stress control's geometry energy and a worse surface, and they do not end more at rest. The stress
control gives exact momentum without a projection, and a better end than a momentum-free force reaches with the same
objective.

\subsection{Comparison with Other Methods}
\label{ch4:sec_comparison}
We compare with nine methods of five kinds, each run from the same sphere toward the same six targets and read by the same
probes: the F-control of Chapter~\ref{CH3_SHAPEMORPHING}~\citep{11088224}; the neural implicit morphs of
\citet{buonomo2025vpnsm}, NISE~\citep{novello2023nise}, ISD~\citep{sang2025isd} and 4Deform~\citep{sang2025fourdeform};
as-rigid-as-possible interpolation~\citep{2000-ARAPShapeInterpolation}; optimal-transport displacement
interpolation~\citep{mccann1997convexity,bonneel2011displacement}; and render-only mesh optimization with Large
Steps~\citep{nicolet2021largesteps} and FlexiCubes~\citep{shen2023flexicubes}. Table~\ref{ch4:tab_methods} and
Appendix~\ref{app:ch4_comparison} give the full reading.

No method's silhouette lies above ours on any mesh, and the clearest separation is in the thin parts: ours leaves about as
many thin target points uncovered as a second sample of the mesh does, while every other method leaves at least half
again as many. What separates the methods most, however, is the way there. ISD fits the end but loses the body on the way;
Large Steps keeps one genus-zero surface and can never open a handle; FlexiCubes changes topology without control;
transport alone and the F-control of Chapter~\ref{CH3_SHAPEMORPHING} keep the mass but tear thin parts into clouds. Ours
stays one body with exact mass and a near-constant volume throughout and reaches every target, including those of higher
genus. Read on their own end meshes, the best mesh methods have sharper surfaces than our particle route, but a perfect
particle sample of the target trails them by a similar amount: that part of the gap is the cost of a particle body at this
resolution. The render-only mesh methods, ARAP, the F-control and NISE are faster than ours; ISD, the volume-preserving
morph and 4Deform are slower.

\subsection{Runtime}
\label{ch4:sec_runtime}
A 300k morph takes ten minutes to half an hour on one GPU (Table~\ref{ch4:tab_runtime}). The run length is set by the stop
rule, and most of it is the slow tail after arrival. That tail is not wasted: the end measures keep improving through its
first half and the body settles, and only its last tenth changes nothing measurable.

\subsection{Thrown and Dropped Bodies}
\label{ch4:sec_throw}
Two spheres are thrown toward each other, collide in the air while half-formed, part, and each becomes a bunny that lands
and tumbles (Figure~\ref{ch4:fig_throw}). The scene uses the multi-body definitions of Appendix~\ref{ch4:sec_scene} and a
stiffer material, with the plastic assimilation off so that the controls hold each body's shape against the impacts. In
the delivered run the bodies meet when the morph is a little under half done, separate cleanly, land, spring off the floor
and tumble, and each ends close to its target in its own frame, with momentum kept through every phase. A single body can
fall as well: a sphere released high above the floor becomes a bunny on the way down, lands, bounces, and tips forward
onto its head, which the body-frame objective allows.

\begin{figure*}[tbp]
    \centering
    \includegraphics[width=\textwidth]{figs/ch4/fig_throw_scene.jpg}
    \caption[{Two thrown bodies, each becoming a bunny}]{\textbf{Two thrown bodies, each becoming a bunny.} Each body has its own grid field and its own objective in its
    own frame, so they part cleanly, fly, land and tumble.}
    \label{ch4:fig_throw}
\end{figure*}

% ============================================================================
\section{Measured Costs and Next Steps}
\label{ch4:sec_costs}
% ============================================================================

\textbf{Rest under the render's authority.} With its weight held after arrival, the render writes what the physics
objective cannot see, but the body is still moving when the run stops, because a held weight does not vanish once the
body has arrived. When such a run's controls are released, the body comes to rest and the colour layout stays where the
render put it, but part of the carved relief is lost; at 40k about a third of it remains. Two parts of the method can take
it back, and they are being separated: the elastic part of the stress the controls built up, and the relaxation of the
outer layer toward the 3D target, which keeps acting after the release and pulls the layer toward a target that is smooth. Taking the written end as the rest shape keeps the lettering, as in the 500k press of
Figure~\ref{ch4:fig_stamp}; whether a run reaches that rest on its own is being measured. The other definition under test
is a render weight that follows the render term's own excess over a measured floor, so that it fades once the
picture-only content is written.
%% SWAP-IN: the settle and rest verdict (job 50) goes here when it lands.
%% SWAP-IN: the stiffness result (job 51), if it bears on the rest, goes here when it lands.

\textbf{Narrow gaps.} The cubic kernel averages the grid velocity over about two cells, so a gap of the target narrower
than that cannot open until late and tears into beads when it does. An MPM cell that follows $N$ opens the teapot's handle
and sharpens the boundaries, but leaves most of the dragon's loops closed, fills the early surface with holes and more
than doubles the dragon's cost. A finer cell is therefore an option of the pipeline, to be chosen after the same comparison
on more meshes.

\textbf{A target given only by pictures.} When the 3D target is the visual hull of the training silhouettes and only the
render terms see the true mesh, the full method sharpens the held-out outline but does not carve the concavities that only
the pictures hold. The information is in the pictures, but only a smooth change of the whole surface lowers the loss; a
multiview shape-from-shading descent finds the concave regions in place, and the next step carves them from a hull that
contains the object at the render's resolution.


Further measured costs, the transient floaters of the early arrival, the ambient drag's effect on angular momentum and the
long, slow tail of the delivered frames, are described in Appendix~\ref{app:ch4_costs}.

% ============================================================================
\section{Chapter Summary}
\label{ch4:sec_summary}
% ============================================================================

This chapter presented a render-guided morphing pipeline built on the differentiable MLS-MPM of
Chapter~\ref{CH3_SHAPEMORPHING}. The body is a volume-conserving solid, optimized in settled windows that are scored after
the controls are released. The stress increment $\delta\mathbf{F}_c$ is its bulk control: it acts through the grid, keeps
the momentum and carries the transport. The layer offset $u$ is its surface control: it moves the outer particles along
their normals, below the grid and free of rigid modes. The render loss asks for what a viewer sees, the silhouette and
shading of the body's own exterior from 18 views, and it reaches $u$ directly and $\delta\mathbf{F}_c$ through the adjoint
of the driven steps; both paths carry it.

Under the default rule of its weight, the render redirects material while the body arrives, then only polishes rims
because its weight follows the vanishing physics gradient, and it leaves the end consistently but slightly closer to the
exact target, with part of its lead on its own term coming from a smoother arrangement rather than a different shape.
With its weight held after arrival, it carves the words SIGGRAPH and VANDERBILT, which exist only in the target pictures,
into a plaque whose 3D target is smooth, and the words are in the material: the plaque prints them into a soft slab,
where a decal on its twin prints nothing. In the same way, it moves red material to the half of a bunny that the pictures
paint red, a layout the physics objective cannot see. Both come at the cost of the body's rest, and a render weight that
keeps the authority and the rest is the next step.

The morph keeps exact mass and, up to the ambient drag, momentum; it ends at rest, opens handles a sphere does not have,
runs unchanged on a wide range of targets and shape pairs, and takes part in contact and gravity with other bodies.
Against nine other methods, no method's silhouette lies above ours on any mesh, and only ours keeps thin parts at the
sampling floor.

% Dissertation build: the stretch is now set for the whole document in the preamble (2026-10-22d); keep it.
\setlength{\emergencystretch}{1.5em}
